\section{hello}

<div class="section">
    <p class="sectionTitle">ប្រធាន</p>
    <p>ចូរសរសេរចំនួនកុំផ្លិចខាងក្រោមជារាង$a+bi$</p>
    <ol class="arabic">
        <li>$-\left( \dfrac{3}{2}+\dfrac{5}{2}i\right)+\left( \dfrac{5}{3}+\dfrac{11}{3}i\right)$</li>
        <li>$\sqrt{-48}-\sqrt{-27}$</li>
        <li>$\left( -1+\sqrt{-3}\right)^2$</li>
        <li>$\left( 3+\sqrt{-5}\right)\left( 7-\sqrt{-10}\right)$</li>
        <li>$\dfrac{2i}{2+i}+\dfrac{5}{2-i}$</li>
        <li>$\dfrac{1+i}{i}-\dfrac{3}{4-i}$</li>
    </ol>
</div>
<div class="section">
    <p class="muolFont">ដំណោះស្រាយ</p>
    <p>សរសេរចំនួនកុំផ្លិចខាងក្រោមជារាង$a+bi$</p>
    <ol class="arabic">
        <li>$-\left( \dfrac{3}{2}+\dfrac{5}{2}i\right)+\left( \dfrac{5}{3}+\dfrac{11}{3}i\right)$<br>
            គេបាន
            $$\begin{aligned}
            &-\left( \dfrac{3}{2}+\dfrac{5}{2}i\right)+\left( \dfrac{5}{3}+\dfrac{11}{3}i\right)\\
            &\quad= -\dfrac{3}{2} - \dfrac{5}{2}i + \dfrac{5}{3} + \dfrac{11}{3}i \\
            &\quad= \left(-\dfrac{3}{2} + \dfrac{5}{3}\right) + \left(-\dfrac{5}{2} + \dfrac{11}{3}\right)i \\
            &\quad= \left(-\dfrac{9}{6} + \dfrac{10}{6}\right) + \left(-\dfrac{15}{6} + \dfrac{22}{6}\right)i \\
            &\quad= \dfrac{1}{6} + \dfrac{7}{6}i
            \end{aligned}$$
            ដូចនេះ
            $$\boxed{-\left( \dfrac{3}{2}+\dfrac{5}{2}i\right)+\left( \dfrac{5}{3}+\dfrac{11}{3}i\right)= \dfrac{1}{6} + \dfrac{7}{6}i}$$
        </li>
        <li>$\sqrt{-48}-\sqrt{-27}$
            <br>
            គេបាន
            $$
            \begin{aligned}
            &\sqrt{-48} - \sqrt{-27} \\
            &\quad= \sqrt{48}i - \sqrt{27}i \\
            &\quad= 4\sqrt{3}i - 3\sqrt{3}i \\
            &\quad= \sqrt{3}i
            \end{aligned}
            $$
            ដូចនេះ$\boxed{\sqrt{-48} - \sqrt{-27} = \sqrt{3}i}$
        </li>
        <li>$\left( -1+\sqrt{-3}\right)^2$
            <br>
            គេបាន
            $$
            \begin{aligned}
            &\left( -1+\sqrt{-3}\right)^2 \\
            &= \left( -1 + \sqrt{3}i \right)^2 \\
            &= \left( 1 - \sqrt{3}i \right)^2 \\
            &= 1 - 2\sqrt{3}i + 3i^2 \\
            &= 1 - 2\sqrt{3}i + 3(-1) \\
            &= 1 - 2\sqrt{3}i - 3 \\
            &= -2 - 2\sqrt{3}i
            \end{aligned}
            $$
            ដូចនេះ $\boxed{\left( -1+\sqrt{-3}\right)^2 = -2 - 2\sqrt{3}i}$
        </li>
        
        <li>$\left( 3+\sqrt{-5}\right)\left( 7-\sqrt{-10}\right)$
            <br>
            គេបាន
            $$
            \begin{aligned}
            &\left( 3+\sqrt{-5}\right)\left( 7-\sqrt{-10}\right) \\
            &= (3 + \sqrt{5}i)(7 - \sqrt{10}i) \\
            &= 21 - 3\sqrt{10}i + 7\sqrt{5}i - (\sqrt{5} \times \sqrt{10})i^2 \\
            &= 21 - 3\sqrt{10}i + 7\sqrt{5}i - (\sqrt{50})(-1) \\
            &= 21 - 3\sqrt{10}i + 7\sqrt{5}i + \sqrt{50} \\
            &= (21 + 5\sqrt{2}) + (7\sqrt{5} - 3\sqrt{10})i
            \end{aligned}
            $$
            ដូចនេះ $$\boxed{\left( 3+\sqrt{-5}\right)\left( 7-\sqrt{-10}\right) = (21 + 5\sqrt{2}) + (7\sqrt{5} - 3\sqrt{10})i}$$
        </li>
        <li>$\dfrac{2i}{2+i}+\dfrac{5}{2-i}$
            <br>
            គេបាន
            $$
            \begin{aligned}
            &\dfrac{2i}{2+i}+\dfrac{5}{2-i} \\
            &= \dfrac{2i}{2+i} + \dfrac{5}{2-i} \\
            &= \dfrac{2i(2-i)}{(2+i)(2-i)} + \dfrac{5(2+i)}{(2-i)(2+i)} \\
            &= \dfrac{2i(2-i) + 5(2+i)}{(2+i)(2-i)} \\
            &= \dfrac{(4i - 2i^2) + (10 + 5i)}{4 - i^2} \\
            &= \dfrac{4i - 2(-1) + 10 + 5i}{4 - (-1)} \\
            &= \dfrac{4i + 2 + 10 + 5i}{5} \\
            &= \dfrac{(12) + (9i)}{5} \\
            &= \dfrac{12}{5} + \dfrac{9}{5}i
            \end{aligned}
            $$
            ដូចនេះ $$\boxed{\dfrac{2i}{2+i}+\dfrac{5}{2-i} = \dfrac{12}{5} + \dfrac{9}{5}i}$$
        </li>
        <li>$\dfrac{1+i}{i}-\dfrac{3}{4-i}$
            <br>
            គេបាន
            $$
            \begin{aligned}
            &\dfrac{1+i}{i} - \dfrac{3}{4-i} \\
            &= \dfrac{(1+i)\cdot (-i)}{i \cdot (-i)} - \dfrac{3(4+i)}{(4-i)(4+i)} \\
            &= \dfrac{(1+i)(-i)}{-i^2} - \dfrac{3(4+i)}{16 - (-1)} \\
            &= \dfrac{-i - i^2}{1} - \dfrac{3(4+i)}{17} \\
            &= -i - (-1) - \dfrac{12 + 3i}{17} \\
            &= 1 - i - \dfrac{12 + 3i}{17} \\
            &= \dfrac{17}{17} - \dfrac{17i}{17} - \dfrac{12}{17} - \dfrac{3i}{17} \\
            &= \dfrac{17 - 12}{17} - \dfrac{17i + 3i}{17} \\
            &= \dfrac{5}{17} - \dfrac{20}{17}i
            \end{aligned}
            $$
            ដូចនេះ $$\boxed{\dfrac{1+i}{i}-\dfrac{3}{4-i} = \dfrac{5}{17} - \dfrac{20}{17}i}$$
        </li>
    </ol>
</div>
\section{hello}
<div class="section">
    <p class="sectionTitle">ឧទាហរណ៍</p>
    <p>ដោះស្រាយប្រព័ន្ធវិសមីការ$\begin{cases}x^2-6x+5\le 0\\ x^2-9x+14\lt 0 \end{cases}$</p>
    <p class="muolFont">ដំណោះស្រាយ</p>
    <p>ដោះស្រាយប្រព័ន្ធវិសមីការ$\begin{cases}x^2-6x+5\le 0\\ x^2-9x+14\lt 0 \end{cases}$<br>
    តាង $f(x)=x^2-6x+5$ <br>
    តាង $g(x)=x^2-9x+14$
    $$\begin{aligned}
        f(x)&=0\Rightarrow x^2-6x+5=0\\
        &\text{ដោយ}\,a+b+c=0\,\text{នោះ}\\
        &x_1=1,\,x_2=\frac{c}{a}=5\\
        g(x)&=0\Rightarrow x^2-9x+14=0\\
        &\Delta=b^2-4ac\\
        &\quad=(-9)^2-4(1)(14)\\
        &\quad=81-56=25\\
        &x_1=\frac{-b-\sqrt{\Delta}}{2a}\\
        &\quad=\frac{-(-9)-\sqrt{25}}{2(1)}\\
        &\quad=\frac{9-5}{2}=2\\
        &x_2=\frac{-b+\sqrt{\Delta}}{2a}\\
        &\quad=\frac{-(-9)+\sqrt{25}}{2(1)}\\
        &\quad=\frac{9+5}{2}=7
    \end{aligned}$$
    តារាងសញ្ញា
    $$\small\begin{array}{c|lcccccccccr}
     x & -\infty & & 1 & & 2 & & 5 & & 7 & & +\infty \\\hline
       & & &\smash{\large\mid}& & \smash{\large\mid} & & \smash{\large\mid} & & \smash{\large\mid} & & \\
     f(x)& & + & 0 & - & \smash{\large\mid} & - & 0 & + & \smash{\large\mid} & + & \\
       & & &\smash{\large\mid}& & \smash{\large\mid} & & \smash{\large\mid} & & \smash{\large\mid} & & \\\hline
       & & &\smash{\large\mid}& & \smash{\large\mid} & & \smash{\large\mid} & & \smash{\large\mid} & & \\
    g(x)& & + &\smash{\large\mid}& + & 0 & - &\smash{\large\mid}&- &0& + & \\
         & & &\smash{\large\mid}& & \smash{\large\mid} & & \smash{\large\mid} & & \smash{\large\mid} & & \\
    \end{array}$$
    តាមតារាង
    $$\begin{aligned}
        &\begin{cases}f(x)\le 0\\ g(x)\lt 0 \end{cases}\\
        \Leftrightarrow& \begin{cases}x\in [1,5]\\ x\in (2,7) \end{cases}\\
        \Leftrightarrow &x\in [1,5]\cap (2,7)\\
        \Leftrightarrow &x\in (2,5]
    \end{aligned}$$
    ដូចនេះ ប្រព័ន្ធវិសមីការមានចម្លើយ$x\in (2,5]$។
    </p>
    
</div>
<div class="section">
    <p class="sectionTitle">ឧទាហរណ៍</p>
    <p>ដោះស្រាយប្រព័ន្ធវិសមីការ$\begin{cases}-x^2+3x+4\ge 0\\ \left(2x^2+3\right)\left(x^2-3x\right) \gt 0 \end{cases}$</p>
    <p class="muolFont">ដំណោះស្រាយ</p>
    <p>ដោះស្រាយប្រព័ន្ធវិសមីការ$\begin{cases}-x^2+3x+4\ge 0\\ \left(2x^2+3\right)\left(x^2-3x\right) \gt 0 \end{cases}$<br>
    ដោយ$2x^2+3 \gt 0$នោះ$\left(2x^2+3\right)\left(x^2-3x\right) \gt 0$$\Leftrightarrow x^2-3x\gt 0$<br>
    ប្រព័ន្ធវិសមីការទៅជា$\begin{cases}-x^2+3x+4\ge 0\\ x^2-3x \gt 0 \end{cases}$<br>
    តាង $f(x)=-x^2+3x+4$ <br>
    តាង $g(x)=x^2-3x$
    $$\begin{aligned}
        f(x)&=0\Rightarrow -x^2+3x+4=0\\
        &\text{ដោយ}\,a-b+c=0\,\text{នោះ}\\
        &x_1=-1,\,x_2=-\frac{c}{a}=4\\
        g(x)&=0\Rightarrow x^2-3x=0\\
        &x(x-3)=0\\
        &x=0\,\text{ឬ}\,x-3=0\\
        &x=0\,\text{ឬ}\,x=3
    \end{aligned}$$
    តារាងសញ្ញា
    $$\small\begin{array}{c|lcccccccccr}
     x & -\infty & & -1 & & 0 & & 3 & & 4 & & +\infty \\\hline
       & & &\smash{\large\mid}& & \smash{\large\mid} & & \smash{\large\mid} & & \smash{\large\mid} & & \\
     f(x)& & - & 0 & + & \smash{\large\mid} & + &\smash{\large\mid}& + &0& - & \\
       & & &\smash{\large\mid}& & \smash{\large\mid} & & \smash{\large\mid} & & \smash{\large\mid} & & \\\hline
       & & &\smash{\large\mid}& & \smash{\large\mid} & & \smash{\large\mid} & & \smash{\large\mid} & & \\
    g(x)& & + &\smash{\large\mid}& + & 0 & - &0&+&\smash{\large\mid}& + & \\
         & & &\smash{\large\mid}& & \smash{\large\mid} & & \smash{\large\mid} & & \smash{\large\mid} & & \\
    \end{array}$$
    តាមតារាង
    $$\begin{aligned}
        &\begin{cases}f(x)\ge 0\\ g(x)\gt 0 \end{cases}\\
        \Leftrightarrow& \begin{cases}x\in [-1,4]\\ x\in (-\infty,0)\cup (3,+\infty) \end{cases}\\
        \Leftrightarrow &x\in [-1,0)\cup (3,4]
    \end{aligned}$$
    ដូចនេះ ប្រព័ន្ធវិសមីការមានចម្លើយ$x\in [-1,0)\cup (3,4]$។
    </p>
</div>

